\IfFileExists{stacks-project.cls}{%
\documentclass{stacks-project}
}{%
\documentclass{amsart}
}

% For dealing with references we use the comment environment
\usepackage{verbatim}
\newenvironment{reference}{\comment}{\endcomment}
%\newenvironment{reference}{}{}
\newenvironment{slogan}{\comment}{\endcomment}
\newenvironment{history}{\comment}{\endcomment}

% For commutative diagrams we use Xy-pic
\usepackage[all]{xy}

% We use 2cell for 2-commutative diagrams.
\xyoption{2cell}
\UseAllTwocells

% We use multicol for the list of chapters between chapters
\usepackage{multicol}

% This is generally recommended for better output
\usepackage{lmodern}
\usepackage[T1]{fontenc}

% For cross-file-references
\usepackage{xr-hyper}

% Package for hypertext links:
\usepackage{hyperref}

% For any local file, say "hello.tex" you want to link to please
% use \externaldocument[hello-]{hello}
\externaldocument[introduction-]{introduction}
\externaldocument[conventions-]{conventions}
\externaldocument[sets-]{sets}
\externaldocument[categories-]{categories}
\externaldocument[topology-]{topology}
\externaldocument[sheaves-]{sheaves}
\externaldocument[sites-]{sites}
\externaldocument[stacks-]{stacks}
\externaldocument[fields-]{fields}
\externaldocument[algebra-]{algebra}
\externaldocument[brauer-]{brauer}
\externaldocument[homology-]{homology}
\externaldocument[derived-]{derived}
\externaldocument[simplicial-]{simplicial}
\externaldocument[more-algebra-]{more-algebra}
\externaldocument[smoothing-]{smoothing}
\externaldocument[modules-]{modules}
\externaldocument[sites-modules-]{sites-modules}
\externaldocument[injectives-]{injectives}
\externaldocument[cohomology-]{cohomology}
\externaldocument[sites-cohomology-]{sites-cohomology}
\externaldocument[dga-]{dga}
\externaldocument[dpa-]{dpa}
\externaldocument[sdga-]{sdga}
\externaldocument[hypercovering-]{hypercovering}
\externaldocument[schemes-]{schemes}
\externaldocument[constructions-]{constructions}
\externaldocument[properties-]{properties}
\externaldocument[morphisms-]{morphisms}
\externaldocument[coherent-]{coherent}
\externaldocument[divisors-]{divisors}
\externaldocument[limits-]{limits}
\externaldocument[varieties-]{varieties}
\externaldocument[topologies-]{topologies}
\externaldocument[descent-]{descent}
\externaldocument[perfect-]{perfect}
\externaldocument[more-morphisms-]{more-morphisms}
\externaldocument[flat-]{flat}
\externaldocument[groupoids-]{groupoids}
\externaldocument[more-groupoids-]{more-groupoids}
\externaldocument[etale-]{etale}
\externaldocument[chow-]{chow}
\externaldocument[intersection-]{intersection}
\externaldocument[pic-]{pic}
\externaldocument[weil-]{weil}
\externaldocument[adequate-]{adequate}
\externaldocument[dualizing-]{dualizing}
\externaldocument[duality-]{duality}
\externaldocument[discriminant-]{discriminant}
\externaldocument[derham-]{derham}
\externaldocument[local-cohomology-]{local-cohomology}
\externaldocument[algebraization-]{algebraization}
\externaldocument[curves-]{curves}
\externaldocument[resolve-]{resolve}
\externaldocument[models-]{models}
\externaldocument[functors-]{functors}
\externaldocument[equiv-]{equiv}
\externaldocument[pione-]{pione}
\externaldocument[etale-cohomology-]{etale-cohomology}
\externaldocument[proetale-]{proetale}
\externaldocument[relative-cycles-]{relative-cycles}
\externaldocument[more-etale-]{more-etale}
\externaldocument[trace-]{trace}
\externaldocument[crystalline-]{crystalline}
\externaldocument[spaces-]{spaces}
\externaldocument[spaces-properties-]{spaces-properties}
\externaldocument[spaces-morphisms-]{spaces-morphisms}
\externaldocument[decent-spaces-]{decent-spaces}
\externaldocument[spaces-cohomology-]{spaces-cohomology}
\externaldocument[spaces-limits-]{spaces-limits}
\externaldocument[spaces-divisors-]{spaces-divisors}
\externaldocument[spaces-over-fields-]{spaces-over-fields}
\externaldocument[spaces-topologies-]{spaces-topologies}
\externaldocument[spaces-descent-]{spaces-descent}
\externaldocument[spaces-perfect-]{spaces-perfect}
\externaldocument[spaces-more-morphisms-]{spaces-more-morphisms}
\externaldocument[spaces-flat-]{spaces-flat}
\externaldocument[spaces-groupoids-]{spaces-groupoids}
\externaldocument[spaces-more-groupoids-]{spaces-more-groupoids}
\externaldocument[bootstrap-]{bootstrap}
\externaldocument[spaces-pushouts-]{spaces-pushouts}
\externaldocument[spaces-chow-]{spaces-chow}
\externaldocument[groupoids-quotients-]{groupoids-quotients}
\externaldocument[spaces-more-cohomology-]{spaces-more-cohomology}
\externaldocument[spaces-simplicial-]{spaces-simplicial}
\externaldocument[spaces-duality-]{spaces-duality}
\externaldocument[formal-spaces-]{formal-spaces}
\externaldocument[restricted-]{restricted}
\externaldocument[spaces-resolve-]{spaces-resolve}
\externaldocument[formal-defos-]{formal-defos}
\externaldocument[defos-]{defos}
\externaldocument[cotangent-]{cotangent}
\externaldocument[examples-defos-]{examples-defos}
\externaldocument[algebraic-]{algebraic}
\externaldocument[examples-stacks-]{examples-stacks}
\externaldocument[stacks-sheaves-]{stacks-sheaves}
\externaldocument[criteria-]{criteria}
\externaldocument[artin-]{artin}
\externaldocument[quot-]{quot}
\externaldocument[stacks-properties-]{stacks-properties}
\externaldocument[stacks-morphisms-]{stacks-morphisms}
\externaldocument[stacks-limits-]{stacks-limits}
\externaldocument[stacks-cohomology-]{stacks-cohomology}
\externaldocument[stacks-perfect-]{stacks-perfect}
\externaldocument[stacks-introduction-]{stacks-introduction}
\externaldocument[stacks-more-morphisms-]{stacks-more-morphisms}
\externaldocument[stacks-geometry-]{stacks-geometry}
\externaldocument[moduli-]{moduli}
\externaldocument[moduli-curves-]{moduli-curves}
\externaldocument[examples-]{examples}
\externaldocument[exercises-]{exercises}
\externaldocument[guide-]{guide}
\externaldocument[desirables-]{desirables}
\externaldocument[coding-]{coding}
\externaldocument[obsolete-]{obsolete}
\externaldocument[fdl-]{fdl}
\externaldocument[index-]{index}

% Theorem environments.
%
\theoremstyle{plain}
\newtheorem{theorem}[subsection]{Theorem}
\newtheorem{proposition}[subsection]{Proposition}
\newtheorem{lemma}[subsection]{Lemma}

\theoremstyle{definition}
\newtheorem{definition}[subsection]{Definition}
\newtheorem{example}[subsection]{Example}
\newtheorem{exercise}[subsection]{Exercise}
\newtheorem{situation}[subsection]{Situation}

\theoremstyle{remark}
\newtheorem{remark}[subsection]{Remark}
\newtheorem{remarks}[subsection]{Remarks}

\numberwithin{equation}{subsection}

% Macros
%
\def\lim{\mathop{\mathrm{lim}}\nolimits}
\def\colim{\mathop{\mathrm{colim}}\nolimits}
\def\Spec{\mathop{\mathrm{Spec}}}
\def\Hom{\mathop{\mathrm{Hom}}\nolimits}
\def\Ext{\mathop{\mathrm{Ext}}\nolimits}
\def\SheafHom{\mathop{\mathcal{H}\!\mathit{om}}\nolimits}
\def\SheafExt{\mathop{\mathcal{E}\!\mathit{xt}}\nolimits}
\def\Sch{\mathit{Sch}}
\def\Mor{\mathop{\mathrm{Mor}}\nolimits}
\def\Ob{\mathop{\mathrm{Ob}}\nolimits}
\def\Sh{\mathop{\mathit{Sh}}\nolimits}
\def\NL{\mathop{N\!L}\nolimits}
\def\CH{\mathop{\mathrm{CH}}\nolimits}
\def\proetale{{pro\text{-}\acute{e}tale}}
\def\etale{{\acute{e}tale}}
\def\QCoh{\mathit{QCoh}}
\def\Ker{\mathop{\mathrm{Ker}}}
\def\Im{\mathop{\mathrm{Im}}}
\def\Coker{\mathop{\mathrm{Coker}}}
\def\Coim{\mathop{\mathrm{Coim}}}

% Boxtimes
%
\DeclareMathSymbol{\boxtimes}{\mathbin}{AMSa}{"02}

%
% Macros for moduli stacks/spaces
%
\def\QCohstack{\mathcal{QC}\!\mathit{oh}}
\def\Cohstack{\mathcal{C}\!\mathit{oh}}
\def\Spacesstack{\mathcal{S}\!\mathit{paces}}
\def\Quotfunctor{\mathrm{Quot}}
\def\Hilbfunctor{\mathrm{Hilb}}
\def\Curvesstack{\mathcal{C}\!\mathit{urves}}
\def\Polarizedstack{\mathcal{P}\!\mathit{olarized}}
\def\Complexesstack{\mathcal{C}\!\mathit{omplexes}}
% \Pic is the operator that assigns to X its picard group, usage \Pic(X)
% \Picardstack_{X/B} denotes the Picard stack of X over B
% \Picardfunctor_{X/B} denotes the Picard functor of X over B
\def\Pic{\mathop{\mathrm{Pic}}\nolimits}
\def\Picardstack{\mathcal{P}\!\mathit{ic}}
\def\Picardfunctor{\mathrm{Pic}}
\def\Deformationcategory{\mathcal{D}\!\mathit{ef}}
\usepackage{cleveref}
\crefname{originalsection}{original section}{original sections}
\Crefname{originalsection}{Original section}{Original sections}
\makeatletter
\newlabel{section-relative-fields}{{0H4E}{source}{}{sourcepointer-1}{}}
\newlabel{section-relative-fields@cref}{{[originalsection][0][]0H4E}{[1][0][]source}{}{}{}}
\makeatother
\crefname{originallemma}{original lemma}{original lemmas}
\Crefname{originallemma}{Original lemma}{Original lemmas}
\makeatletter
\newlabel{lemma-multiplicities-field-extension}{{0H4F}{source}{}{sourcepointer-2}{}}
\newlabel{lemma-multiplicities-field-extension@cref}{{[originallemma][0][]0H4F}{[1][0][]source}{}{}{}}
\makeatother
\crefname{originallemma}{original lemma}{original lemmas}
\Crefname{originallemma}{Original lemma}{Original lemmas}
\makeatletter
\newlabel{lemma-coequalizer-dim-r}{{0H4W}{source}{}{sourcepointer-3}{}}
\newlabel{lemma-coequalizer-dim-r@cref}{{[originallemma][0][]0H4W}{[1][0][]source}{}{}{}}
\makeatother
\title{Descent of families of r-cycles along morphisms with separable residue field extensions}
\author{The Stacks Project Authors}
\date{stacks-project:relative-cycles.tex:local-snapshot-20261012; GNU FDL 1.2 or later}
\begin{document}
\maketitle
\setcounter{subsection}{1}
\noindent\textit{Editorial note.} All delivered mathematical text is the author's own, copied from the pinned Stacks Project chapter file relative-cycles.tex; no statement or proof step was rewritten, reordered or removed. The counted span is the author environment and its own proof environment. References inside the span are handed to the assembler as navigation stubs that carry the canonical tag of the cited tag only; whether each of them is load-bearing is recorded on the stub itself.\par\medskip
\noindent\textit{Allowed standard prerequisites (editorial).}\begin{itemize}
\item Injectivity of base change for cycles over fields: for a locally algebraic scheme X over a field k, the base change map Z\_r(X) to Z\_r(X\_\{k'\}) on r-cycles is always injective (relative-cycles chapter, Section section-relative-fields, canonical tag 0H4E). cited at proof lines 848-850. Yields part (1): the base change map on families of r-cycles over X/S is injective whenever S' to S is surjective.
\item Positivity of multiplicities in field extensions: the multiplicity m\_\{Z', Z\_K\} of an irreducible component Z' of the base change of an integral locally algebraic k-scheme Z by a field extension K/k is one or a power of the characteristic of k, in particular a positive integer (relative-cycles chapter, Lemma lemma-multiplicities-field-extension, canonical tag 0H4F). cited at proof lines 880-881. The coefficients m(x', g) in the base change operation on coefficient functions are positive integers.
\item Existence and coequalizer property of lifts of points of transdegree r under a surjective base change: every x of transdegree r admits a lift x' of transdegree r, and any two such lifts are both projections of a common x'' of transdegree r in the double base change (relative-cycles chapter, Lemma lemma-coequalizer-dim-r, canonical tag 0H4W). cited at proof lines 899-903. Together with the displayed multiplicity identity and the footnote construction, gives a unique rational-valued function alpha with g\textasciicircum{}*alpha = alpha'.
\end{itemize}\medskip
\ifdefined\hypertarget\hypertarget{sourcepointer-1}{}\fi
\noindent\textit{Source reference.} Original section 0H4E: sources/stacks-chow/tags.txt, source lines 21349--21349. Cited at proof lines 848-850: part (1) rests on the injectivity of the base change map established there.\par\medskip
\ifdefined\hypertarget\hypertarget{sourcepointer-2}{}\fi
\noindent\textit{Source reference.} Original lemma 0H4F: sources/stacks-chow/tags.txt, source lines 21350--21350. Cited at proof lines 880-881: source of the positivity of the coefficients m(x', g) in the base change operation on coefficient functions.\par\medskip
\ifdefined\hypertarget\hypertarget{sourcepointer-3}{}\fi
\noindent\textit{Source reference.} Original lemma 0H4W: sources/stacks-chow/tags.txt, source lines 21366--21366. Cited at proof lines 899-903: with the footnote construction yields the unique rational-valued function alpha with g\textasciicircum{}*alpha = alpha'.\par\medskip
\begin{lemma}
\label{lemma-descend-family}
Let $f : X \to S$ be a morphism of schemes which is locally of finite type.
Let $r \geq 0$ be an integer. Let $g : S' \to S$ be a morphism of
schemes and $X' = S' \times_S X$. Assume that for every $s \in S$ there
exists a point $s' \in S'$ with $g(s') = s$ and such that
$\kappa(s')/\kappa(s)$ is a separable extension of fields. Then
\begin{enumerate}
\item For families $\alpha_1$ and $\alpha_2$ of $r$-cycles on fibres of $X/S$
if $g^*\alpha_1 = g^*\alpha_2$, then $\alpha_1 = \alpha_2$.
\item Given a family $\alpha'$ of $r$-cycles on fibres of $X'/S'$ if
$\text{pr}_1^*\alpha' = \text{pr}_2^*\alpha'$ as families of
$r$-cycles on fibres of $(S' \times_S S') \times_S X / (S' \times_S S')$,
then there is a unique family $\alpha$ of $r$-cycles on fibres of $X/S$
such that $g^*\alpha = \alpha'$.
\end{enumerate}
\end{lemma}

\begin{proof}
Part (1) follows from the injectivity of the base change map discussed
in Section \ref{section-relative-fields}. (This argument works as
long as $S' \to S$ is surjective.)

\medskip\noindent
Let $\alpha'$ be as in (2). Denote
$\alpha'' = \text{pr}_1^*\alpha' = \text{pr}_2^*\alpha'$
the common value.

\medskip\noindent
Let $(X/S)^{(r)}$ be the set of $x \in X$ with
$\text{trdeg}_{\kappa(f(x))}(\kappa(x)) = r$
and similarly define $(X'/S')^{(r)}$ and $(X''/S'')^{(r)}$
Taking coefficients, we may think of $\alpha'$ and $\alpha''$ as functions
$\alpha' : (X'/S')^{(r)} \to \mathbf{Z}$ and
$\alpha'' : (X''/S'')^{(r)} \to \mathbf{Z}$.
Given a function
$$
\varphi : (X/S)^{(r)} \to \mathbf{Z}
$$
we define $g^*\varphi : (X'/S')^{(r)} \to \mathbf{Z}$ by analogy
with our base change operation. Namely, say $x' \in (X'/S')^{(r)}$
maps to $x \in X$, $s' \in S'$, and $s \in Z$.
Denote $Z' \subset X'_{s'}$ and $Z \subset X_s$ the integral
closed subschemes with generic points $x'$ and $x$. Note
that $\dim(Z') = r$. If $\dim(Z) < r$, then we set $(g^*\varphi)(x') = 0$.
If $\dim(Z) = r$, then $Z'$ is an irreducible component of $Z_{s'}$ and
hence has a multiplicity $m_{Z', Z_{s'}}$. Call this $m(x', g)$.
Then we define
$$
(g^*\varphi)(x') = m(x', g) \varphi(x)
$$
Note that the coefficients $m(x', g)$ are always positive integers
(see for example Lemma \ref{lemma-multiplicities-field-extension}).
We similarly have base change maps
$$
\text{pr}_1^*, \text{pr}_2^* :
\text{Map}((X'/S')^{(r)}, \mathbf{Z})
\longrightarrow
\text{Map}((X''/S'')^{(r)}, \mathbf{Z})
$$
It follows from the associativity of base change that we have
$\text{pr}_1^* \circ g^* = \text{pr}_2^* \circ g^*$ (small detail
omitted). To be explicitly, in terms of the maps of sets this
equality just means that for $x'' \in (X''/S'')^{(r)}$ we have
$$
m(x'', \text{pr}_1) m(\text{pr}_1(x''), g) =
m(x'', \text{pr}_2) m(\text{pr}_2(x''), g)
$$
provided that $\text{pr}_1(x'')$ and $\text{pr}_2(x'')$ are
in $(X''/S'')^{(r)}$.
By Lemma \ref{lemma-coequalizer-dim-r} and an elementary
argument\footnote{Given $x \in (X/S)^{(r)}$ pick $x' \in (X'/S')^{(r)}$
mapping to $x$ and set $\alpha(x) = \alpha'(x')/m(x', g)$. This
is well defined by the formula and the lemma.}
using the previous displayed equation,
it follows that there exists a unique map
$$
\alpha : (X/S)^{(r)} \to \mathbf{Q}
$$
such that $g^*\alpha = \alpha'$. To finish the proof it suffices
to show that $\alpha$ has integer values (small detail omitted: one
needs to see that $\alpha$ determines a locally finite sum on
each fibre which follows from the corresponding fact for $\alpha'$).
Given any $x \in (X/S)^{(r)}$ with image $s \in S$
we can pick a point $s' \in S'$ such that $\kappa(s')/\kappa(s)$
is separable. Then we may choose $x' \in (X'/S')^{(r)}$ mapping
to $s$ and $x$ and we see that $m(x', g) = 1$ because
$Z_{s'}$ is reduced in this case. Whence $\alpha(x) = \alpha'(x')$
is an integer.
\end{proof}

\end{document}
